%============================================================
%  Bd. 21 - Rechtskürzbare Halbgruppen
%============================================================

\documentclass[main.tex]{subfiles}

\ifSubfilesClassLoaded{
  \BandSetupStandalone{B21}
}{
}

\title{Bd. 21 - Rechtskürzbare Halbgruppen}
\author{Martin Kunze}
\date{}
\setcounter{file}{21}

\begin{document}

\title{Bd. 21 - Rechtskürzbare Halbgruppen}
\author{Martin Kunze}
\date{}
\hypersetup{pageanchor=false}
\maketitle
\hypersetup{pageanchor=true}
\setcounter{tocdepth}{3}
\tableofcontents
\setcounter{file}{21}
\renewcommand{\FormulaBandID}{21}

\chapter{Einführung}

Dieser Band untersucht Halbgruppen mit rechtskürzbarer Operation.

\chapter{Rechtskürzbare Halbgruppen}

\section{Definition der Rechtskürzbarkeit}

\FormulaDefDeltaK[Rechtskürzbare Operation]{%
  \mathsf{RK\ddot{u}rz}\!\left(A,\star\right)%
}{RightCancellativeDef}{%
  \DeltaRow{Mengen}{A\dsep a\dsep b\dsep c}
  \DeltaRow{Rechtskürzung}{%
    a,b,c\in A\dsep b\star a=c\star a\vdash b=c}
  \DeltaRow{\textbf{Neues Symbol}}{}
  \DeltaPrem{Rechtskürzbarkeit}{\mathsf{RK\ddot{u}rz}\!\left(A,\star\right)}
}

\FormulaAxiomDeltaK[Zugriff auf die Rechtskürzung]{%
  \mathsf{RK\ddot{u}rz}\!\left(A,\star\right)\dsep
  a,b,c\in A\dsep b\star a=c\star a
  \vdash b=c%
}{RightCancellationAxiom}{%
  \DeltaRow{Mengen}{A\dsep a\dsep b\dsep c}
}

\FormulaDefDeltaK[Begriff der rechtskürzbaren Halbgruppe]{%
  \RightCancellativeSemigroupAlg{A}{\star}%
}{RightCancellativeSemigroupDef}{%
  \DeltaRow{Mengen}{A}
  \DeltaRow{\textbf{Axiome}}{}
  \DeltaPrem{Halbgruppe}{\SemigroupAlg{A}{\star}}
  \DeltaPrem{Rechtskürzbarkeit}{%
    \mathsf{RK\ddot{u}rz}\!\left(A,\star\right)}
  \DeltaRow{\textbf{Neues Symbol}}{}
  \DeltaPrem{Rechtskürzbare Halbgruppe}{%
    \RightCancellativeSemigroupAlg{A}{\star}}
}

\FormulaAxiomDeltaK[Zugriff auf die Halbgruppe]{%
  \RightCancellativeSemigroupAlg{A}{\star}
  \vdash\SemigroupAlg{A}{\star}%
}{RightCancellativeSemigroupBaseAxiom}{%
  \DeltaRow{Mengen}{A}
}

\FormulaAxiomDeltaK[Zugriff auf die Rechtskürzbarkeit]{%
  \RightCancellativeSemigroupAlg{A}{\star}
  \vdash\mathsf{RK\ddot{u}rz}\!\left(A,\star\right)%
}{RightCancellativeSemigroupPropertyAxiom}{%
  \DeltaRow{Mengen}{A}
}

\FormulaAxiomDeltaK[Direkte Rechtskürzung]{%
  \RightCancellativeSemigroupAlg{A}{\star}\dsep
  a,b,c\in A\dsep b\star a=c\star a
  \vdash b=c%
}{RightCancellativeSemigroupCancellationAxiom}{%
  \DeltaRow{Mengen}{A\dsep a\dsep b\dsep c}
}

\section{Beispiele}

\FormulaThmDeltaK[Die Addition ist rechtskürzbar]{%
  \mathsf{RK\ddot{u}rz}\!\left(\mathbb{N},+\right)%
}{NaturalAdditionRightCancellative}{%
  \DeltaRow{Mengen}{\mathbb{N}\dsep a\dsep b\dsep c}
}
\begin{tabproof}
  \proofstep{1}{a\in\mathbb{N}}{\rA}
  \proofstep{2}{b\in\mathbb{N}}{\rA}
  \proofstep{3}{c\in\mathbb{N}}{\rA}
  \proofstep{4}{b+a=c+a}{\rA}
  \proofstep{1,2,3,4}{b=c}{%
    \FormulaRefAuto{PeanoAddRightCancellation}[thm]{2,3,1,4}}
  \proofstep{}{\mathsf{RK\ddot{u}rz}\!\left(\mathbb{N},+\right)}{%
    \FormulaRefAuto{RightCancellativeDef}[def]{5}}
\end{tabproof}

\FormulaThmDeltaK[Die Addition bildet eine rechtskürzbare Halbgruppe]{%
  \RightCancellativeSemigroupAlg{\mathbb{N}}{+}%
}{NaturalAdditionRightCancellativeSemigroup}{%
  \DeltaRow{Mengen}{\mathbb{N}}
}
\begin{tabproof}
  \proofstep{}{\SemigroupAlg{\mathbb{N}}{+}}{%
    \FormulaRefAuto{NaturalAdditionSemigroup}[thm]}
  \proofstep{}{\mathsf{RK\ddot{u}rz}\!\left(\mathbb{N},+\right)}{%
    \FormulaRefAuto{NaturalAdditionRightCancellative}[thm]}
  \proofstep{}{\RightCancellativeSemigroupAlg{\mathbb{N}}{+}}{%
    \FormulaRefAuto{RightCancellativeSemigroupDef}[def]{1,2}}
\end{tabproof}

\paragraph{Einseitiges Beispiel.}
Für $x\star_\ell y\coloneqq x$ ist jede nichtleere Menge $A$
rechtskürzbar. Besitzt $A$ mindestens zwei Elemente, ist diese Halbgruppe
nicht linkskürzbar.

\section{Zusammenhang mit der Linkskürzbarkeit}

\FormulaThmDeltaK[Rechtskürzungsinstanz aus Linkskürzung und Kommutativität]{%
  \begin{aligned}[t]
    &\CommutativeSemigroupAlg{A}{\star}\dsep
      \mathsf{LK\ddot{u}rz}\!\left(A,\star\right)\dsep
      a,b,c\in A\dsep b\star a=c\star a\vdash{}\\[-2pt]
    &b=c
  \end{aligned}
}{CommutativeLeftCancellationImpliesRightCancellation}{%
  \DeltaRow{Mengen}{A\dsep a\dsep b\dsep c}
}
\begin{tabproofwide}
  \proofstepwidestar[1]{\CommutativeSemigroupAlg{A}{\star}}{\rA}
  \proofstepwidestar[2]{\mathsf{LK\ddot{u}rz}\!\left(A,\star\right)}{\rA}
  \proofstepwidestar[3]{a\in A}{\rA}
  \proofstepwidestar[4]{b\in A}{\rA}
  \proofstepwidestar[5]{c\in A}{\rA}

  \proofstepwide[1,3,4]{a\star b}{=}{b\star a}{%
    \FormulaRefAuto{CommutativeSemigroupCommutativityAxiom}[ax]{1,3,4}}
  \proofstepwide[7]{}{=}{c\star a}{\rA}
  \proofstepwide[1,3,5]{}{=}{a\star c}{%
    \FormulaRefAuto{CommutativeSemigroupCommutativityAxiom}[ax]{1,5,3}}
  \proofstepwidestar[1,3,4,5,7]{a\star b=a\star c}{%
    \rChain{6,8}}

  \proofstepwidestar[1,2,3,4,5,7]{b=c}{%
    \FormulaRefAuto{LeftCancellationAxiom}[ax]{2,3,4,5,9}}
\end{tabproofwide}

\FormulaThmDeltaK[Kommutative linkskürzbare Halbgruppen sind rechtskürzbar]{%
  \CommutativeSemigroupAlg{A}{\star}\dsep
  \mathsf{LK\ddot{u}rz}\!\left(A,\star\right)
  \vdash\mathsf{RK\ddot{u}rz}\!\left(A,\star\right)%
}{CommutativeLeftCancellativeImpliesRightCancellative}{%
  \DeltaRow{Mengen}{A\dsep a\dsep b\dsep c}
}
\begin{tabproof}
  \proofstep{1}{\CommutativeSemigroupAlg{A}{\star}}{\rA}
  \proofstep{2}{\mathsf{LK\ddot{u}rz}\!\left(A,\star\right)}{\rA}
  \proofstep{3}{a\in A}{\rA}
  \proofstep{4}{b\in A}{\rA}
  \proofstep{5}{c\in A}{\rA}
  \proofstep{6}{b\star a=c\star a}{\rA}
  \proofstep{1,2,3,4,5,6}{b=c}{%
    \FormulaRefAuto{CommutativeLeftCancellationImpliesRightCancellation}[thm]{%
      1,2,3,4,5,6}}
  \proofstep{1,2}{\mathsf{RK\ddot{u}rz}\!\left(A,\star\right)}{%
    \FormulaRefAuto{RightCancellativeDef}[def]{7}}
\end{tabproof}

\FormulaThmDeltaK[Kommutative linkskürzbare Halbgruppen sind
rechtskürzbare Halbgruppen]{%
  \CommutativeSemigroupAlg{A}{\star}\dsep
  \LeftCancellativeSemigroupAlg{A}{\star}
  \vdash\RightCancellativeSemigroupAlg{A}{\star}%
}{CommutativeLeftCancellativeSemigroupImpliesRightCancellativeSemigroup}{%
  \DeltaRow{Mengen}{A}
}
\begin{tabproof}
  \proofstep{1}{\CommutativeSemigroupAlg{A}{\star}}{\rA}
  \proofstep{2}{\LeftCancellativeSemigroupAlg{A}{\star}}{\rA}
  \proofstep{2}{\SemigroupAlg{A}{\star}}{%
    \FormulaRefAuto{LeftCancellativeSemigroupBaseAxiom}[ax]{2}}
  \proofstep{2}{\mathsf{LK\ddot{u}rz}\!\left(A,\star\right)}{%
    \FormulaRefAuto{LeftCancellativeSemigroupPropertyAxiom}[ax]{2}}
  \proofstep{1,2}{\mathsf{RK\ddot{u}rz}\!\left(A,\star\right)}{%
    \FormulaRefAuto{CommutativeLeftCancellativeImpliesRightCancellative}[thm]{1,4}}
  \proofstep{1,2}{\RightCancellativeSemigroupAlg{A}{\star}}{%
    \FormulaRefAuto{RightCancellativeSemigroupDef}[def]{3,5}}
\end{tabproof}

\section{Morphismen}

\FormulaDefDeltaK[Homomorphismus rechtskürzbarer Halbgruppen]{%
  \RightCancellativeSemigroupHom{f}{A,\star}{B,\diamond}%
}{RightCancellativeSemigroupHomDef}{%
  \DeltaRow{Mengen}{A\dsep B}
  \DeltaRow{Funktionen}{f\dsep\star\dsep\diamond}
  \DeltaRow{\textbf{Axiome}}{}
  \DeltaPrem{Quellstruktur}{\RightCancellativeSemigroupAlg{A}{\star}}
  \DeltaPrem{Zielstruktur}{\RightCancellativeSemigroupAlg{B}{\diamond}}
  \DeltaPrem{Halbgruppenhomomorphismus}{%
    \SemigroupHom{f}{A,\star}{B,\diamond}}
  \DeltaRow{\textbf{Neues Symbol}}{}
  \DeltaPrem{Homomorphismus}{%
    \RightCancellativeSemigroupHom{f}{A,\star}{B,\diamond}}
}

\FormulaAxiomDeltaK[Quellstruktur]{%
  \RightCancellativeSemigroupHom{f}{A,\star}{B,\diamond}
  \vdash\RightCancellativeSemigroupAlg{A}{\star}%
}{RightCancellativeSemigroupHomSourceAxiom}{%
  \DeltaRow{Mengen}{A\dsep B}
  \DeltaRow{Funktionen}{f\dsep\star\dsep\diamond}
}

\FormulaAxiomDeltaK[Zielstruktur]{%
  \RightCancellativeSemigroupHom{f}{A,\star}{B,\diamond}
  \vdash\RightCancellativeSemigroupAlg{B}{\diamond}%
}{RightCancellativeSemigroupHomTargetAxiom}{%
  \DeltaRow{Mengen}{A\dsep B}
  \DeltaRow{Funktionen}{f\dsep\star\dsep\diamond}
}

\FormulaAxiomDeltaK[Zugriff auf den Halbgruppenhomomorphismus]{%
  \RightCancellativeSemigroupHom{f}{A,\star}{B,\diamond}
  \vdash\SemigroupHom{f}{A,\star}{B,\diamond}%
}{RightCancellativeSemigroupHomBaseHomAxiom}{%
  \DeltaRow{Mengen}{A\dsep B}
  \DeltaRow{Funktionen}{f\dsep\star\dsep\diamond}
}

\FormulaDefDeltaK[Isomorphismus rechtskürzbarer Halbgruppen]{%
  \RightCancellativeSemigroupIso{f}{A,\star}{B,\diamond}%
}{RightCancellativeSemigroupIsoDef}{%
  \DeltaRow{Mengen}{A\dsep B}
  \DeltaRow{Funktionen}{f\dsep\star\dsep\diamond}
  \DeltaRow{\textbf{Axiome}}{}
  \DeltaPrem{Homomorphismus}{%
    \RightCancellativeSemigroupHom{f}{A,\star}{B,\diamond}}
  \DeltaPrem{Bijektivität}{f\colon A\bij B}
  \DeltaRow{\textbf{Neues Symbol}}{}
  \DeltaPrem{Isomorphismus}{%
    \RightCancellativeSemigroupIso{f}{A,\star}{B,\diamond}}
}

\FormulaAxiomDeltaK[Homomorphismuszugriff]{%
  \RightCancellativeSemigroupIso{f}{A,\star}{B,\diamond}
  \vdash\RightCancellativeSemigroupHom{f}{A,\star}{B,\diamond}%
}{RightCancellativeSemigroupIsoHomAxiom}{%
  \DeltaRow{Mengen}{A\dsep B}
  \DeltaRow{Funktionen}{f\dsep\star\dsep\diamond}
}

\FormulaAxiomDeltaK[Bijektivitätszugriff]{%
  \RightCancellativeSemigroupIso{f}{A,\star}{B,\diamond}
  \vdash f\colon A\bij B%
}{RightCancellativeSemigroupIsoBijectiveAxiom}{%
  \DeltaRow{Mengen}{A\dsep B}
  \DeltaRow{Funktionen}{f\dsep\star\dsep\diamond}
}

\end{document}
